\subsection{補完実験とseed 0--19の統合集計}

2モデル補完実験（seed 0--9）は1,000件を記録し，L2・L3成功376件，失敗624件だった．QPE seed 2の自己デバッグ最終roundに残った1件は，状態ベクトルを保存せずResultから取得した生成コード側の誤用である．保存済み例外を根拠にexec / interface\_mismatchへ読み込み時に再分類し，生データ・生成コード・評価値は保持した．全件のsanity・automation再検証が成功し，API障害・timeout・未分類エラーは0件となった．

検証済み3実験の1,500＋2,500＋1,000件を統合し，5モデル×10タスク×20 seed×5条件＝5,000件を得た．各モデル1,000件，各モデル・条件200件で，欠落・重複・予定外記録は0件．モデルdigest，生成設定，プロンプト，評価環境の一致を確認した．ただしschema 2.0と2.1の診断文処理差を含むため，単一手順の介入効果を示すものではない．出典別内訳を保持した混合手順の記述的集計として報告する．

\begin{table}[t]\centering

\caption{統合集計のL2・L3成功件数（各セル200件）}\label{tab:pooled}

\begin{tabular}{l|rrrrr}\hline

モデル & B & R & F & D & P \\\hline

Q4 & 35 & 67 & 109 & 118 & 25 \\

G4 & 133 & 167 & 183 & 176 & 63 \\

Q9 & 116 & 144 & 157 & 158 & 58 \\

C7 & 132 & 148 & 158 & 155 & 129 \\

G3 & 2 & 2 & 2 & 2 & 18 \\

\hline

合計 & 418 & 528 & 609 & 609 & 293 \\\hline

\end{tabular}\end{table}

全条件のL2・L3成功は2,457件（49.14\%）で一致した．Bの初期L3成功418件は修正型の各条件で修正せず終了した．初期未達582件を分母とするL3救出はR 110件，F 191件，D 191件である．最終L3退行は修正型では0件だった．Pは別プロンプトの独立生成で，救出85件と退行210件を分けて扱う．

\begin{table}[t]\centering

\caption{統合baseline失敗類型別のL2救出件数}\label{tab:pooledcatalog}

\begin{tabular}{l|r|rrr}\hline

失敗類型 & 分母 & R & F & D \\\hline

import & 271 & 24 & 55 & 65 \\

回路構築 & 106 & 24 & 57 & 57 \\

インターフェース & 32 & 8 & 14 & 13 \\

構文 & 68 & 21 & 25 & 18 \\

出力不一致 & 102 & 33 & 40 & 38 \\

量子ビット数 & 3 & 0 & 0 & 0 \\

\hline

合計 & 582 & 110 & 191 & 191 \\\hline

\end{tabular}\end{table}

importではD，構文と出力不一致ではFの救出件数が観測最多だった．回路構築ではFとDが同数である．モデル・難度・seedブロックの構成と診断文の差が含まれるため，この表だけから最適介入や統計的有意差を断定しない．

難度別のL2達成は易は654/1000件，中は880/2000件，難は923/2000件であった．難度区分は本ハーネスのタスク分類であり，普遍的な難度尺度ではない．
